\section{The settlement window across the year}\label{sec:window}

\subsection{Activity}

Table~\ref{tab:activity} reports how the settlement window on each type of session differs
from the same security's window on ordinary sessions. On an ordinary session the window
accounts for about \ActWindowshareLiquidNormal{}\% of the session's volume in the liquid
group, and similar shares in the other two.

\begin{table}[htbp]
\centering
\small
\caption{Activity in the settlement window by type of session. Each entry is the coefficient
on an indicator for the session type in a regression with security fixed effects, estimated
separately for each group, and measures the difference from ordinary sessions of the same
security. The last column is the difference between the securities with derivatives and the
placebo group, estimated with security and session fixed effects. Standard errors clustered
by session in parentheses. $^{*}$, $^{**}$ and $^{***}$ denote significance at the ten, five
and one per cent levels.}
\label{tab:activity}
\resizebox{\textwidth}{!}{\input{generated/tables/activity}}
\end{table}

On monthly expiries, trading is drawn into the window in securities with derivatives and not
in the others. The window's share of session volume rises by \ActWindowshareLiquidMonthlyexpiry{}
percentage points in the liquid group and by \ActWindowshareIlliquidMonthlyexpiry{} points
in the illiquid group, while the change in the placebo group, \ActWindowsharePlaceboMonthlyexpiry{}
points, is not distinguishable from zero. The difference between the derivatives groups and
the placebo group is \ActWindowshareDidMonthlyexpiry{} points, with a standard error of
\ActWindowshareDidMonthlyexpirySe{}. The calendar controls do not reproduce this pattern. On weekly
index expiries the window's share of volume does not change in any group. On month ends it
rises sharply in all three groups, by \ActWindowshareLiquidMonthend{} points in the liquid
group and by \ActWindowsharePlaceboMonthend{} points in the placebo group, which reflects the
closing demand of index and portfolio rebalancing; the additional rise in the securities with
derivatives, \ActWindowshareDidMonthend{} points, is smaller than on monthly expiries.

The settlement price also moves further from the pre-window price on monthly expiries. For
the securities with derivatives the absolute move increases by
\ActAbsdriftDerivativeMonthlyexpiry{} basis points from an average of
\ActAbsdriftDerivativeNormal{}, and the final five minutes of the window end
\ActAbsterminalDerivativeMonthlyexpiry{} basis points further from the settlement price.
Neither increase is specific to the derivatives groups. The placebo group shows increases of
similar size, and the differences between the two, \ActAbsdriftDidMonthlyexpiry{} and
\ActAbsterminalDidMonthlyexpiry{} basis points, are small and insignificant. Figure~\ref{fig:daytypes}
shows the same comparison in levels.

\begin{figure}[htbp]
  \centering
  \includegraphics[width=\textwidth]{day_types.pdf}
  \caption{Share of session volume traded in the settlement window, and the absolute move
    from the pre-window price to the settlement price, by group and type of session. Means
    over security-sessions.}
  \label{fig:daytypes}
\end{figure}

\subsection{Is the settlement price reversed?}

Table~\ref{tab:reversal} estimates equation~\eqref{eq:reversal}. On ordinary sessions the move
over the window is not reversed the next morning: the coefficient for the securities with
derivatives is $\RevDerivativeDrift{}$, with a standard error of \RevDerivativeDriftSe{}. On
monthly expiries a part of the move is reversed. The additional coefficient is
$\RevDerivativeDriftxmonthlyexpiry{}$ (standard error \RevDerivativeDriftxmonthlyexpirySe{}),
so that about a quarter of the window's move is undone by the next morning, and the
randomization $p$-value is \RevDerivativeDriftxmonthlyexpiryRip{}. The corresponding
coefficients for weekly index expiries and month ends are smaller and not significant.

\begin{table}[htbp]
\centering
\small
\caption{Reversal of the settlement price. Regressions of the return from the settlement price
to the volume-weighted average price of the next session's first thirty minutes on the move
over the window and its interactions with the session types, with security and session fixed
effects. The last column is the difference between the securities with derivatives and the
placebo group. Standard errors clustered by session in parentheses. The randomization
$p$-value compares the monthly-expiry coefficient with its distribution over \RiDraws{}
samples in which the expiries are replaced by one ordinary session from each month.}
\label{tab:reversal}
\resizebox{\textwidth}{!}{\input{generated/tables/reversal}}
\end{table}

The reversal is not, however, confined to securities with an expiring contract. In the placebo
group the additional coefficient on monthly expiries is $\RevPlaceboDriftxmonthlyexpiry{}$
(randomization $p$-value \RevPlaceboDriftxmonthlyexpiryRip{}), and the difference between the
two sets of securities, $\RevDidDriftxmonthlyexpiry{}$, is close to zero (randomization
$p$-value \RevDidDriftxmonthlyexpiryRip{}). A settlement price displaced by trading aimed at
the settlement should be reversed where there is a settlement to influence and not elsewhere.
The pattern found here is a property of the monthly expiry session as a whole, shared by
securities on which nothing settles.

Figure~\ref{fig:path} shows the same result as a path. For each security and session the
prices through the window and at the next morning's open are measured relative to the average
of all securities on that session, and oriented in the direction of the window's move. On
ordinary sessions the move is kept overnight; on monthly expiries part of it is given back, to
a similar extent in all three groups.

\begin{figure}[htbp]
  \centering
  \includegraphics[width=\textwidth]{window_path.pdf}
  \caption{Average path of the price through the settlement window and at the next morning's
    open, relative to the pre-window price and to the average of all securities on the same
    session, oriented so that the window's move is positive.}
  \label{fig:path}
\end{figure}

\subsection{Whose flow moves the settlement price?}

Table~\ref{tab:flow} relates the window's move and the next morning's return to the net
aggressor flow in the window. Across the securities with derivatives, net aggressive buying of
one per cent of median session volume raises the settlement price by
\FlowDriftAllDerivativeFlow{} basis points on an ordinary session, and the effect on monthly
expiries is not different.

The participant categories differ sharply. Aggressive flow of custodians and of other clients
moves the settlement price in the direction of the flow, by \FlowDriftParticipantDerivativeFlowcustodian{}
and \FlowDriftParticipantDerivativeFlowncnp{} basis points per per cent of volume. Aggressive
flow of proprietary traders is associated with a move in the opposite direction,
$\FlowDriftParticipantDerivativeFlowproprietary{}$ basis points, and is followed by a
next-morning return of \FlowNextrevParticipantDerivativeFlowproprietary{} basis points in its
own favour. Proprietary traders thus trade aggressively against the window's move and are
rewarded overnight, which is the behaviour of liquidity supply by arbitrageurs rather than of
price pressure. None of these relations changes on monthly expiries: no participant category's
window flow has a larger price effect, or is followed by a larger reversal, on the days on
which the settlement is struck.

\begin{table}[htbp]
\centering
\small
\caption{Net aggressor flow in the window, the window's move and the next morning's return.
Flow is the volume bought by aggressive buyers less the volume sold by aggressive sellers, in
per cent of the security's median session volume, in total and by the aggressor's category.
Regressions include security and session fixed effects and the session-type indicators and
their interactions with flow; only the flow terms and their interaction with the monthly
expiry are shown. Standard errors clustered by session in parentheses.}
\label{tab:flow}
\resizebox{\textwidth}{!}{\input{generated/tables/flow}}
\end{table}
